\label{abstract}

\textbf{Bối cảnh:} Quá trình đô thị hóa nhanh chóng và sự xuống cấp của hạ tầng xây dựng đã và đang làm gia tăng nguy cơ tiềm ẩn về cháy nổ tại các khu dân cư và khu công nghiệp. Trong bối cảnh đó, lính cứu hỏa thường xuyên phải đối mặt với các hiểm nguy nghiêm trọng như nhiệt độ cực cao, khí độc, không gian hẹp và nguy cơ sập đổ công trình trong khi thi hành nhiệm vụ. Các phương tiện liên lạc truyền thống như bộ đàm hai chiều chỉ cung cấp thông tin phân mảnh và thiếu khả năng giám sát trạng thái từng cá nhân theo thời gian thực, dẫn đến những hệ quả nghiêm trọng khi không xác định được vị trí hoặc điều kiện của lính cứu hỏa — đặc biệt trong các tình huống bất tỉnh hoặc bất động. Nhận dạng hoạt động con người (Human Activity Recognition -- HAR) kết hợp với dữ liệu vị trí là một hướng giải pháp tiềm năng, cung cấp thông tin khách quan và liên tục để hỗ trợ ra quyết định kịp thời cho trung tâm chỉ huy.

\textbf{Mục tiêu:} Đề tài hướng tới xây dựng hệ thống nhận dạng hoạt động lính cứu hỏa theo thời gian thực (Real-time Firefighter Activity Recognition -- RFAR) có khả năng phân loại chính xác mười một hoạt động đặc thù trong các tình huống cứu hộ khẩn cấp. Hệ thống được thiết kế dựa trên một thiết bị đeo tích hợp cảm biến gia tốc kế gắn tại vùng thắt lưng, đảm bảo tính nhỏ gọn, tiết kiệm năng lượng và khả năng triển khai thực tế trên vi điều khiển có tài nguyên hạn chế.

\textbf{Phương pháp:} Nghiên cứu xây dựng một tập dữ liệu riêng mô phỏng các hoạt động đặc thù của lính cứu hỏa với sự tham gia của 20 tình nguyện viên nam trong các tình huống diễn tập thực tế tại cơ sở đào tạo phòng cháy chữa cháy. Một quy trình xử lý tín hiệu mới được đề xuất, kết hợp kỹ thuật cửa sổ trượt 3 giây với thuật toán phát hiện chỉ số bất thường (Abnormality Index -- AIx) nhằm xử lý hiệu quả các hoạt động chuyển tiếp và sự kiện ngã. Các đặc trưng miền thời gian tối ưu được lựa chọn để giảm độ phức tạp tính toán, phù hợp với yêu cầu triển khai trên vi điều khiển. Thuật toán Random Forest (RF) được hiệu chỉnh siêu tham số chặt chẽ nhằm đảm bảo kích thước mô hình dưới 1~MB trong khi vẫn duy trì hiệu suất nhận dạng cao.

\textbf{Kết quả bước đầu:} Mô hình RF tối ưu với 17 cây quyết định và độ sâu tối đa 14 đạt kích thước 0,92~MB, độ chính xác 99,3\% và điểm F1-score 96,1\% trên tập dữ liệu riêng. Trong các thực nghiệm thời gian thực tại cơ sở đào tạo phòng cháy chữa cháy Việt Nam, hệ thống đạt điểm F1-score 96,2\% với độ trễ suy luận xấp xỉ 9,6~ms trên vi điều khiển ESP32. Khi đánh giá trên ba tập dữ liệu công khai, hệ thống đạt F1-score 96,6\% trên SFDLA và MobiFall, và 78,2\% trên UniMiB-SHAR. Hệ thống phát hiện ngã đạt độ chính xác 100\% khi phân biệt ngã (FA) với các hoạt động thông thường.

\textbf{Kết luận:} Đề tài đã hoàn thành các mục tiêu cốt lõi của giai đoạn nghiên cứu hiện tại, bao gồm: xây dựng tập dữ liệu chuyên biệt, phát triển quy trình xử lý tín hiệu mới, tối ưu hóa và triển khai mô hình học máy trên thiết bị nhúng, đồng thời xác nhận tính khả thi qua thực nghiệm thực tế. Trong các giai đoạn tiếp theo, nghiên cứu sẽ tập trung vào tối ưu hóa thiết kế phần cứng, mở rộng đối tượng thu thập dữ liệu và tích hợp hệ thống định vị trong nhà cùng cơ chế cảnh báo khẩn cấp tự động.